\scanpage{302}
\uline{Exemple 2} \deleted{Supposons} $X_1=X_2=P^r$, et $\varphi_1,\varphi_2$ sont deux \deleted{\unsure{applications}} immersions \unsure{correspondants} à deux \deleted{\unsure{classes de}} classes $\unsure{\mathcal L_i}$ de $H^1(Y,\uline{\mathcal O}_Y^*)$ \deleted{\unsure{telles que}} qui sont \uline{indépendantes sur $\mathbb Z$} : On peut par exemple supposer $r=2$, $Y$ une courbe elliptique. Alors \deleted{\unsure{un}} tt faisceau inversible sur $X_1\amalg_Y X_2$ est \uline{trivial}, car il induit un faisceau trivial sur $X_1$ \deleted{$\unsure{\times}$}\inserted{et} $X_2$, et l'isom. \unsure{requis} de leurs \deleted{\unsure{restrictions}} images inverses à $Y$, moyennant identifications \uline{convenables} des $\unsure{\uline{L}_i|\mathcal F_i}\unsure{\simeq|:}\uline{\mathcal O}_{X_i}$ $(i=1,2)$ est l'identité. En particulier, $X_1\amalg_Y X_2$ \uline{n'est pas} \uline{projectif}, bien que son normalisé $X_1\amalg_{\deleted{\unsure{Y}}}X_2=\deleted{\unsure{P}}\inserted{P^r\amalg P^r}$ soit.

Ici encore, \deleted{\unsure{faisant opérer} $\mathbb Z/2$} on peut faire un passage au quotient plus compliqué,
\figure{\begin{tikzpicture}\draw[->] (0.4,1.4)--(0.1,1);\draw (1.5,0)--(1.5,0.35)--(2.6,0.4);\end{tikzpicture}}
\mspage{301}
